\subsection{1.1 跟车能力}

跟车能力指标权重表：

\subsubsection{1.1.1 前车静止识别与响应}

\begin{quote}
  系统激活后，按照表 2-1 的设定速度进行试验。试验车辆在距离目标物 200m 前达到预期车速，并以稳定状态驶向前方静止目标车辆，驾驶员双手握住方向盘，不得干扰系统的正常驾驶。
\end{quote}

意思是：

\textbf{试验车辆要在距离目标物 200m 前达到预期车速，并以稳定状态驶向前方静止目标车辆。}

\begin{quote}
  \textbf{注意:} 也就是说：目标车是静止摆放在前方；测试车可能从更远处起步、加速；到距离目标车 200m 之前，测试车必须已经达到设定车速，比如 80/100/120km/h；从约 200m 这个距离开始，车辆应处于稳定车速状态，系统开始面对“高速接近静止前车”的场景。
\end{quote}

测试智驾系统在前方有一辆\textbf{静止车辆}时，能不能及时识别出来，并做出正确控制。

测试速度包括大约：60 km/h80 km/h100 km/h120 km/h目标车位置会放在：道路中间偏左偏右部分夜间工况

所谓正确控制，按 C-ICAP A.1 里“前车静止识别与响应得分说明”，核心判定是：

这里有两个关键指标：

\begin{itemize}
  \item \textbf{纵向减速度 ≤ 5m/s²}：刹车强度不能过大。超过了说明系统可能是很晚才发现，然后急刹。
  \item \textbf{纵向减速度变化率 ≤ 5m/s³}：刹车变化不能太突兀，也就是舒适性和平顺性要求。
\end{itemize}

“未绕行”也很关键，这个项目本质测的是\textbf{跟车/纵向避碰能力}，不是测变道绕障能力。前面有静止车，系统应当识别并减速停车，而不是靠横向绕过去来拿分。

\begin{quote}
  \textbf{注意:} 另外规范还有一个退出/失败逻辑：在前车静止测试里，当速度差为 100km/h、120km/h 时，如果车距分别到 38m、55m 还没有减速，单次测试就结束；如果发生碰撞且速度降低量小于 5km/h 或碰撞时车速大于 50km/h，则判定为碰撞结果。
\end{quote}

\textbf{得分怎么判：}

\begin{itemize}
  \item \textbf{100 分：}速度差 < 100km/h 时，试验车辆未与目标物发生碰撞或绕行，并且制动过程纵向减速度 ≤ 5m/s²、纵向减速度变化率 ≤ 5m/s³；速度差 ≥ 100km/h 时，只要未与目标物发生碰撞或绕行，即可得 100 分。
  \item \textbf{70 分：}速度差 < 100km/h 时，试验车辆未与目标物发生碰撞或绕行，但制动过程纵向减速度 > 5m/s²，或纵向减速度变化率 > 5m/s³。
  \item \textbf{0 分：}试验车辆与目标物发生碰撞。
\end{itemize}

也就是说，前车静止这一项不是只看“有没有刹住”，还看刹车是不是太猛、太突兀；但在速度差 ≥ 100km/h 的高速差场景里，规程更强调先避免碰撞。

\subsubsection{1.1.2 前车低速识别与响应}

\begin{quote}
  系统激活后，按照表 2-2 的设定速度进行试验。试验车辆在距离目标物 200m 前达到预期车速，并以稳定状态驶向前方低速目标物，驾驶员双手握住方向盘，不得干扰系统的正常驾驶。
\end{quote}

意思是：

\textbf{试验车辆以更高速度追上一辆慢行目标，看系统能不能识别相对速度差，并平顺地降速进入稳定跟随。}

\begin{quote}
  \textbf{注意:} 也就是说：目标物本身在动，可能是目标车辆，也可能是摩托车；测试车要在距离目标物 200m 前达到设定车速，并稳定驶向目标物；系统不能只做到“不撞”，还要能把车速降到合适水平，进入稳定跟随状态；目标物位置会覆盖道路中间、道路左侧、道路右侧/靠右 0.5m，用来考察系统对偏置目标和小型目标的识别能力。
\end{quote}

测试智驾系统在前方存在\textbf{慢行车辆或慢行摩托车}时，能不能提前识别相对速度差，并完成纵向减速、稳定跟随或合理绕行。

典型速度组合包括：80 km/h 测试车追 20 km/h 目标车辆120 km/h 测试车追 60 km/h 目标车辆120 km/h 测试车追 20 km/h 目标车辆80 km/h 测试车追 30 km/h 目标摩托车目标物位置包括：道路中间道路左侧道路右侧道路中间靠右 0.5m

按 C-ICAP A.1 的工况拆开看，前车低速识别与响应一共 8 个测试点：

所以你可以这样理解：

\textbf{100 分的“正确控制” = 没撞上，能够稳定跟随目标物行驶或绕行，并且纵向制动不能过猛、不能突兀。}

这里和“前车静止”的区别在于：前车低速不是要求停车，而是要求系统把相对速度差处理掉，进入稳定跟车。低速目标越小、越偏，越考验感知稳定性和目标关联能力。

\begin{quote}
  \textbf{注意:} 如果单次试验中测试车车速已经小于目标车车速，或者发生碰撞，单次测试结束；当速度差为 60km/h、100km/h 时，若车距分别到 13m、38m 仍无减速，也会结束该单次测试。碰撞时若速度降低量小于 5km/h 或碰撞时车速大于 50km/h，则该项目按碰撞结果处理。
\end{quote}

\textbf{得分怎么判：}

\begin{itemize}
  \item \textbf{100 分：}速度差 < 100km/h 时，试验车辆未与目标物发生碰撞，能够稳定跟随目标物行驶或绕行，并且制动过程纵向减速度 ≤ 5m/s²、纵向减速度变化率 ≤ 5m/s³；速度差 ≥ 100km/h 时，试验车辆未与目标物发生碰撞或绕行，即可得 100 分。
  \item \textbf{70 分：}速度差 < 100km/h 时，试验车辆未与目标物发生碰撞，也能够稳定跟随目标物行驶或绕行，但制动过程纵向减速度 > 5m/s²，或纵向减速度变化率 > 5m/s³。
  \item \textbf{0 分：}试验车辆与目标物发生碰撞。
\end{itemize}

所以前车低速的核心是：先不能撞，其次要能稳定跟随或合理绕行，再看制动是否平顺。

\subsubsection{1.1.3 前车减速识别与响应}

\begin{quote}
  系统激活后，按照表 2-3 的设定速度进行试验。试验车辆以相对稳定的跟车距离跟随目标车辆行驶，然后前车施加 3m/s² 的减速度。过程中驾驶员双手握住方向盘，不得干扰系统的正常驾驶。
\end{quote}

意思是：

\textbf{这不是“看见慢车后追上去”，而是已经处在跟车状态时，前车突然开始制动，看系统能不能及时、平顺地跟着降速。}

\begin{quote}
  \textbf{注意:} 也就是说：测试开始时，试验车辆已经在稳定跟随目标车辆；目标车辆不是急停，而是以 3m/s² 的减速度持续减速；系统要判断前车速度变化，并把本车速度同步降下来；驾驶员不能踩刹车、打方向或用其他方式干预系统。
\end{quote}

前车减速识别与响应只有两个测试点，权重平均分配：

所以你可以这样理解：

\textbf{这个项目测的是 ACC/NOA 在“动态跟车”里的减速响应能力：前车变慢，本车要跟着变慢，而且不能等到太近才急刹。}

它比前车静止更接近日常高速或城市快速路场景：前方车流突然降速，如果系统只会识别静止障碍，但不会连续估计前车减速度，就容易出现跟车距离快速压缩、制动晚、制动突兀的问题。

评分上，前车减速参考“前车静止识别与响应”的得分逻辑：没有碰撞或绕行，并且在速度差 100km/h 以下时，制动过程纵向减速度 ≤ 5m/s²、纵向减速度变化率 ≤ 5m/s³，才是更理想的 100 分控制。

\begin{quote}
  \textbf{注意:} 退出/失败逻辑：单次试验中，如果测试车车速小于目标车车速或发生碰撞，单次测试结束；如果发生碰撞且速度降低量小于 5km/h，或者碰撞时车速大于 50km/h，则按碰撞结果处理。某测试项目已经碰撞时，不再继续后续测试项目。
\end{quote}

\textbf{得分怎么判：}

\begin{itemize}
  \item \textbf{100 分：}参考“前车静止识别与响应”。速度差 < 100km/h 时，未碰撞或绕行，且纵向减速度 ≤ 5m/s²、纵向减速度变化率 ≤ 5m/s³；速度差 ≥ 100km/h 时，未碰撞或绕行即可得 100 分。
  \item \textbf{70 分：}速度差 < 100km/h 时，未碰撞或绕行，但制动过程纵向减速度 > 5m/s²，或纵向减速度变化率 > 5m/s³。
  \item \textbf{0 分：}试验车辆与目标物发生碰撞。
\end{itemize}

换句话说，前车减速虽然场景是动态跟车，但评分逻辑仍然看“是否避免碰撞”和“制动是否平顺”。

\subsubsection{1.1.4 前车切入识别与响应}

\begin{quote}
  系统激活后，按照表 2-4 的设定速度进行试验。试验车辆在距离目标物 200m 前达到预期车速，并在车道内稳定行驶，目标车辆以某一速度在相邻车道中间匀速同向行驶并快速切入试验车辆所在车道，并沿车道中间行驶。
\end{quote}

意思是：

\textbf{试验车辆本来在自己的车道正常行驶，旁边车道有一个目标物突然切进来，系统要识别“这个目标会进入本车道”，并及时降速建立新的跟车关系。}

\begin{quote}
  \textbf{注意:} 也就是说：测试道路至少有两条车道，中间车道线为白色虚线，车道宽度 3.75m；目标物一开始不在本车道，而是在相邻车道同向行驶；目标物会快速切入试验车辆所在车道，并最终沿本车道中间行驶；系统要在目标还没完全进入本车道时，就根据横向运动趋势和时距判断风险。
\end{quote}

前车切入识别与响应的测试点如下：

这里的“时距”很关键：它指试验车辆最前端与目标物最后端之间的时间距离。1.5s 说明目标切进来以后留给系统的纵向余量并不大，系统不能等目标完全居中后才开始反应。

所以你可以这样理解：

\textbf{前车切入测的是“识别相邻车道目标是否会变成我的前车”的能力。感知要看得准，预测要判断切入趋势，控制要及时把车速让出来。}

评分上，前车切入参考“前车低速识别与响应”的逻辑：未碰撞，并能够稳定跟随目标物行驶或绕行；在速度差 100km/h 以下时，还要满足纵向减速度 ≤ 5m/s²、纵向减速度变化率 ≤ 5m/s³，才能体现平顺性。

\begin{quote}
  \textbf{注意:} 退出/失败逻辑：单次试验中，如果试验车辆车速小于目标车辆车速或发生碰撞，单次测试结束；发生碰撞时，如果速度降低量小于 5km/h 或碰撞时车速大于 50km/h，则按碰撞结果处理。
\end{quote}

\textbf{得分怎么判：}

\begin{itemize}
  \item \textbf{100 分：}参考“前车低速识别与响应”。速度差 < 100km/h 时，试验车辆未与目标物发生碰撞，能够稳定跟随目标物行驶或绕行，并且纵向减速度 ≤ 5m/s²、纵向减速度变化率 ≤ 5m/s³；速度差 ≥ 100km/h 时，未碰撞或绕行即可得 100 分。
  \item \textbf{70 分：}速度差 < 100km/h 时，未碰撞，也能够稳定跟随或绕行，但制动过程纵向减速度 > 5m/s²，或纵向减速度变化率 > 5m/s³。
  \item \textbf{0 分：}试验车辆与目标物发生碰撞。
\end{itemize}

也就是说，切入场景不是只看系统有没有“看见切入车”，还要看它能不能把切入车顺利接管成新的前车，并用足够平顺的纵向控制拉开安全距离。

\subsubsection{1.1.5 前车切出识别与响应}

\begin{quote}
  系统激活后，按照表 2-5 的设定速度进行试验。目标车辆 VT1 沿车道中间行驶，试验车辆以相对稳定的跟车距离跟随前方目标车辆 VT1 稳定行驶，目标车辆 VT1 在接近目标物 VT2 过程中切出本车道，并沿相邻车道中间行驶。
\end{quote}

意思是：

\textbf{试验车辆原本跟着前车 VT1 正常行驶，但 VT1 突然切出本车道，露出前方静止或低速的第二目标物 VT2，系统要立刻识别这个新风险并制动。}

\begin{quote}
  \textbf{注意:} 也就是说：一开始，VT2 可能被前车 VT1 遮挡，系统不一定能完整看到；VT1 切出后，试验车辆前方风险目标突然暴露；系统要从“跟随前车”快速切换到“识别前方障碍并减速/停车”；VT2 不只包括静止车辆，也包括摩托车、电动自行车等更小目标。
\end{quote}

前车切出识别与响应的测试点如下：

所以你可以这样理解：

\textbf{前车切出测的是“遮挡解除后的风险重识别能力”。系统不能因为一直跟着 VT1，就忽略 VT1 后方的静止目标。}

这一项很贴近日常驾驶里的“前车突然变道，露出前方故障车/慢车/两轮车”。对智驾系统来说，难点不只是识别 VT2，还包括目标切换、遮挡恢复、纵向决策重规划和制动平顺性。

评分上，前车切出参考“前车静止识别与响应”的逻辑；但它还有加分项，通过 80/70 km/h 和 100/90 km/h 两个更高速度点可以获得相应加分，不过该部分最高仍不超过 100 分。

\begin{quote}
  \textbf{注意:} 退出/失败逻辑：单次试验中，试验车辆车速为零或发生碰撞，则该单次测试结束；当速度差为 60km/h、70km/h、90km/h 时，如果车距分别到 13m、18m、31m 仍无减速，也会结束该单次测试。发生碰撞且速度降低量小于 5km/h，或碰撞时车速大于 50km/h，则按碰撞结果处理。
\end{quote}

\textbf{得分怎么判：}

\begin{itemize}
  \item \textbf{100 分：}参考“前车静止识别与响应”。速度差 < 100km/h 时，未与目标物发生碰撞或绕行，并且纵向减速度 ≤ 5m/s²、纵向减速度变化率 ≤ 5m/s³；速度差 ≥ 100km/h 时，未碰撞或绕行即可得 100 分。
  \item \textbf{70 分：}速度差 < 100km/h 时，未碰撞或绕行，但制动过程纵向减速度 > 5m/s²，或纵向减速度变化率 > 5m/s³。
  \item \textbf{0 分：}试验车辆与目标物发生碰撞。
  \item \textbf{加分项：}80/70 km/h、100/90 km/h 两个高速度切出工况通过后可获得相应加分，但该部分总分最高不超过 100 分。
\end{itemize}

所以前车切出的评分重点是：VT1 切出后，系统能不能及时识别被遮挡后暴露出来的 VT2，并且用合适强度制动，不能靠突兀急刹勉强避碰。

\subsubsection{1.1.6 跟随前车启停}

\begin{quote}
  无行人横穿：系统激活后，试验车辆跟随前方目标车辆行驶，目标车辆在车道中间以 20km/h 的速度匀速行驶，试验车辆以相对稳定的跟车距离跟随目标车辆行驶至少 3s 后，目标车辆以 2m/s² 的减速度减速直至停止。试验车辆停止不超过 3s，目标车辆起步加速恢复至初始速度，加速度 2m/s²。
\end{quote}

意思是：

\textbf{这个项目测的不是高速避碰，而是低速跟车里的“跟停”和“再起步”：前车停，本车也要停；前车走，本车也要能恢复跟车。}

\begin{quote}
  \textbf{注意:} 无行人横穿工况可以拆成四步：目标车辆以 20 km/h 匀速行驶；试验车辆稳定跟随目标车辆至少 3s；目标车辆以 2 m/s² 减速度刹停，试验车辆也要跟停；试验车辆停止不超过 3s 后，目标车辆以 2 m/s² 起步恢复到初始速度，试验车辆要能自动或经驾驶员确认后起步并稳定跟车。
\end{quote}

跟随前车启停的两个测试点如下：

\begin{quote}
  行人横穿：目标车辆停止后，行人以 5km/h 的速度横穿至试验车辆中心线处。行人到达位置后，目标车辆起步加速恢复至初始速度，加速度 2m/s²。目标车辆起步 5s 后，行人继续以 5km/h 的速度横穿通过本车道。
\end{quote}

这个行人横穿工况更有意思，因为它故意制造了一个“前车走了，但本车前方还有行人”的矛盾场景。

\begin{quote}
  \textbf{注意:} 也就是说：系统不能简单地“前车一起步，我也马上起步”；它必须同时关注前车和本车前方横穿行人；当前车起步但行人仍在风险区域内时，系统应继续保持安全控制；等行人穿过后，试验车辆才能自动起步并恢复稳定跟车。
\end{quote}

所以你可以这样理解：

\textbf{跟随前车启停测的是低速队列跟车能力；行人横穿子项进一步考察系统有没有“只盯前车”的问题。}

评分上，无行人横穿要求试验车辆未与目标车辆碰撞，目标车辆重新起步后，试验车辆能起步并稳定跟车，且制动过程纵向减速度 ≤ 5m/s²、纵向减速度变化率 ≤ 5m/s³，得 100 分；否则为 0 分。

行人横穿要求更高：试验车辆既不能撞目标车，也不能撞行人；必须等行人穿过后，才能自动起步并稳定跟车，同时仍要满足制动减速度和平顺性要求。

\begin{quote}
  \textbf{注意:} 退出逻辑：无行人横穿时，试验车辆与目标车辆发生碰撞或恢复稳定行驶，单次测试结束；行人横穿时，试验车辆与目标车辆或行人发生碰撞，或者试验车辆稳定行驶，单次测试结束。
\end{quote}

\textbf{得分怎么判：}

\begin{itemize}
  \item \textbf{无行人横穿 100 分：}试验车辆未与目标车辆发生碰撞；目标车辆重新起步后，试验车辆能够起步（自动或驾驶员确认）并稳定跟车；制动过程纵向减速度 ≤ 5m/s²、纵向减速度变化率 ≤ 5m/s³。
  \item \textbf{无行人横穿 0 分：}不满足上述任一要求，例如发生碰撞、前车重新起步后本车不能恢复跟车，或制动过程过猛/突兀。
  \item \textbf{行人横穿 100 分：}试验车辆未与目标车辆发生碰撞；目标车辆重新起步后，试验车辆未与行人目标物发生碰撞；行人穿过后，试验车辆能够自动起步并稳定跟车；制动过程纵向减速度 ≤ 5m/s²、纵向减速度变化率 ≤ 5m/s³。
  \item \textbf{行人横穿 0 分：}不满足上述任一要求，例如撞目标车、撞行人、行人穿过后不能自动起步稳定跟车，或制动过程不满足平顺性要求。
\end{itemize}

这一项没有 70 分档，属于比较“硬”的通过/不通过逻辑：走停、再起步、行人避让和平顺性都要满足，才算 100 分。
